% !TeX spellcheck = de_DE
\chapter{Hardware}
\label{sec:hardware}

\section{Entwicklungsansatz und Design-Entscheidungen}

\subsection{Ausgangssituation}

Das Projekt verfolgt das Ziel, einen kostengünstigen, modularen Roboterararm mit sechs Freiheitsgraden zu entwickeln. Die Hauptanforderungen waren:

\begin{itemize}
	\item \textbf{Modulare Architektur}: Teilweise Unabhängige Entwicklung, was dadurch Austausch einzelner Komponenten vereinfacht
	\item \textbf{Additive Fertigung}: Möglichst einfache Fertigungsart um Diy zu sein
	\item \textbf{Open-Source-Prinzip}: Vollständige CAD-Dateien offen verfügbar
	\item \textbf{Kostengünstigkeit}: Minimierung von gekauften Komponenten durch intelligentes Design
\end{itemize}

\subsection{Design-Aufbau}

\subsubsection{Segmentierter Aufbau}

Der Arm wurde in \textbf{6 funktionale Segmente} unterteilt, die jeweils einen Freiheitsgrad realisieren:

\begin{table}[h]
	\centering
	\begin{tabular}{|l|c|l|p{5cm}|}
		\hline
		\textbf{Segment} & \textbf{DOF} & \textbf{Bewegung} & \textbf{Komponenten} \\
		\hline
		Base & -- & Stativ & Strukturrahmen, Befestigungspunkte \\
		\hline
		Shoulder & 1 & Vertikale Rotation (Yaw) & Motor, Getriebe, Hall Sensor \\
		\hline
		Upper Arm & 2 & Pitch \& Roll &  Motor, Getriebe, Hall Sensor \\
		\hline
		Forearm & 1 & Ellbogen-Pitch & Motor \\
		\hline
		Wrist & 2 & Pitch \& Roll & Motor, Übsetzung \\
		\hline
		Gripper & -- & Greifbewegung & 2-Finger-Greifer, Servo \\
		\hline
	\end{tabular}
	\caption{Segmentierung des Roboterarms mit Freiheitsgraden}
	\label{tab:segments}
\end{table}

\subsubsection{Iterative Verbesserungen (V1 bis V6)}

Der Designprozess durchlief mehrere Versionsschritte:

\begin{itemize}
	\item \textbf{V1--V3}: Konzeptphase mit grundlegenden Geometrien und ersten Motorintegrationen
	\item \textbf{V4}: Optimierung der Gelenkgeometrien, Reduktion von Spielräumen, erste Kollisionstests
	\item \textbf{V5}: Produktionsreife Version mit verfeinerten Toleranzen, optimierter Massenverteilung und vollständiger CAD-Dokumentation
	\item \textbf{V6 Aktuell}: Große Änderungen durch Fehlern welche in V5 waren
\end{itemize}

\section{Konstruktive Lösungsansätze}

\subsection{Design choices und Hardware wahl}
Die Antriebselemente wurden nach folgenden Kriterien ausgewählt:
\begin{itemize}
	\item \textbf{Nema 17 und Nema 23 Motoren}: Für präzise Kontrolle durch Micro-Stepping.
	\item \textbf{Getriebe}: In V6 wurden 20:1 Cycloidal-Drive-Getriebe verwendet, da diese kaum merkbares Spiel aufweisen.
\end{itemize}

\subsubsection*{Berechnung der theoretischen und realen Nutzlast}
Um die tatsächliche Nutzlast (Payload) des Roboterarms am Endeffektor zu ermitteln, muss das theoretische maximale Drehmoment dem statischen Eigengewicht der bewegten Struktur gegenübergestellt werden. Der Basis-Motor wird hierbei nicht berücksichtigt, da sein Gewicht nicht gegen die Schwerkraft gehoben werden muss.

\paragraph{Theoretisches Gesamtdrehmoment:}
Der an der Schulterachse verwendete Nema 23 Schrittmotor liefert ein maximales Haltemoment von $M_{Motor} = 2\text{ Nm}$. Durch die 20-fache Untersetzung des Cycloidal-Getriebes ($i = 20$) ergibt sich ein theoretisches Gesamtdrehmoment am Schultergelenk von:
$$M_{gesamt} = M_{Motor} \cdot i = 2\text{ Nm} \cdot 20 = 40\text{ Nm}$$

Bei einer maximalen Auslage (Reach) von $r_{max} = 0,5\text{ m}$ resultiert daraus eine theoretische maximale Hebekraft $F_{max}$, die ohne Berücksichtigung des Eigengewichts senkrecht nach unten wirken darf:
$$F_{max} = \frac{M_{gesamt}}{r_{max}} = \frac{40\text{ Nm}}{0,5\text{ m}} = 80\text{ N}$$

\paragraph{Einfluss von Eigengewicht und Schwerkraft:}
Von diesem theoretischen Maximum muss das Drehmoment $M_{Eigen}$ abgezogen werden, welches durch das Eigengewicht des ausgestreckten Arms entsteht. Das bewegte Gesamtgewicht $m_{bewegt}$ setzt sich aus den 3D-gedruckten Strukturteilen (ca. $1,91\text{ kg}$) und den mitgehobenen Motoren (ein Nema 23 und drei Nema 17, ca. $2,30\text{ kg}$) zusammen.

Für eine statische Worst-Case-Abschätzung nehmen wir eine Gesamtmasse von $m_{bewegt} = 4,21\text{ kg}$ an. Der gemeinsame Massenschwerpunkt wird zur Vereinfachung in der Mitte der Auslage bei $r_{S} = 0,25\text{ m}$ angenommen. Bei der Erdbeschleunigung ($g = 9,81\text{ m/s}^2$) ergibt sich:
$$M_{Eigen} = m_{bewegt} \cdot g \cdot r_{S}$$
$$M_{Eigen} = 4,21\text{ kg} \cdot 9,81\frac{\text{m}}{\text{s}^2} \cdot 0,25\text{ m} \approx 10,33\text{ Nm}$$

\paragraph{Reale statische Nutzlast:}
Das verbleibende Drehmoment, das für die eigentliche Nutzlast am Endeffektor ($M_{Nutz}$) zur Verfügung steht, ist die Differenz aus dem Gesamtdrehmoment und dem Eigengewicht:
$$M_{Nutz} = M_{gesamt} - M_{Eigen} = 40\text{ Nm} - 10,33\text{ Nm} = 29,67\text{ Nm}$$

Damit ergibt sich die maximal hebbare reale Masse $m_{Nutz}$ am Greifer bei voller Auslage von $0,5\text{ m}$:
$$m_{Nutz} = \frac{M_{Nutz}}{r_{max} \cdot g} = \frac{29,67\text{ Nm}}{0,5\text{ m} \cdot 9,81\frac{\text{m}}{\text{s}^2}} \approx 6,05\text{ kg}$$

\textit{Hinweis: Dies stellt eine rein statische Betrachtung dar. In der Praxis verringern der mechanische Wirkungsgrad der gedruckten Getriebe sowie dynamische Trägheitsmomente bei der Beschleunigung die tatsächliche Nutzlast.}

\section{Zukünftige Entwicklung (V6)}

Version 6 wird mit verbesserten Aspekten geplant:

\begin{itemize}
	\item Optimierte Massenverteilung für präzisere Kinematik
	\item Zusätzliche Sensorsupports (position durch Hallsensor)
	\item Verbesserte Zugänglichkeit für Wartung und Kalibrierung
\end{itemize}


